\section*{Chapter \ref{ch:m3:power-markets-trading} --- Power Markets II: Trading}

\begin{solution}{exo:m3:power-markets-trading:1}
\omterm{def:m3:power-markets-trading:capture}{Capture price} $(800 + 800 + 400 + 900)/100 = \qty{29.00}{EUR/MWh}$; baseload
\qty{50.00}{EUR/MWh}; capture rate 58\%.
\end{solution}

\begin{solution}{exo:m3:power-markets-trading:2}
It pays $20 \times 300 = $ EUR~6{,}000. Long by 20~MWh it would have received EUR~6{,}000 under a
single price.
\end{solution}

\begin{solution}{exo:m3:power-markets-trading:3}
Positions are set a day ahead on forecasts that keep changing: wind, solar, demand and plant
outages. Intraday trading lets each party correct its position as its forecast improves,
instead of paying \omterm{def:m3:power-markets-trading:imbalance}{imbalance prices} for the whole error.
\end{solution}

\begin{solution}{exo:m3:power-markets-trading:4}
$50 \times (15 - (-5)) \times 10 = $ EUR~10{,}000.
\end{solution}

\begin{solution}{exo:m3:power-markets-trading:5}
It draws $200/\sqrt{0.88} = 213.2$~MWh at \qty{20}{}, paying EUR~4{,}264, and delivers
$200\sqrt{0.88} = 187.6$~MWh at \qty{120}{}, receiving EUR~22{,}514: EUR~18{,}250.
\end{solution}

\begin{solution}{exo:m3:power-markets-trading:6}
$(60 - 46.07) \times 100\,000 = $ EUR~1.39 million a year.
\end{solution}

\begin{solution}{exo:m3:power-markets-trading:7}
EUR~65{,}968 per MW in 2024 and EUR~74{,}546 in 2025 with one cycle a day; EUR~88{,}688 in 2025
with two.
\end{solution}

\begin{solution}{exo:m3:power-markets-trading:8}
The backtest assumes perfect foresight of the day's prices, a price-taking battery that does
not move the auction, and no degradation, outages or fees; real operators forecast, and their
own bids narrow the spreads they trade. It also ignores the intraday and reserve revenues that
may matter more. Treat it as an upper bound on day-ahead revenue.
\end{solution}

\begin{solution}{pb:m3:power-markets-trading:1}
\textbf{1.} It charges from 14:00 to 16:00 (prices $-2.30$ and $-2.17$) and discharges from 20:00
to 22:00 (140.99 and 175.32): the cheapest pair of hours followed by the dearest pair.
\textbf{2.} It draws 213.2~MWh and delivers 187.6~MWh. \textbf{3.} EUR~30{,}149. \textbf{4.} At a
negative price it is paid to charge; the lowest prices were negative that afternoon.
\textbf{5.} EUR~760 more: the second-best spread that day was small.
\textbf{6.} EUR~65{,}968 in 2024 and EUR~74{,}546 in 2025 per MW. \textbf{7.} More negative and
near-zero midday hours as solar grew, and high evening prices: the daily spread widened.
\textbf{8.} EUR~14{,}142 per MW. \textbf{9.} Spreads are widest when solar is high and evenings
are tight (spring and summer) and in volatile winter weeks. \textbf{10.} Without losses it would
have earned EUR~32{,}078 on 18 June: the 12\% loss costs about EUR~1{,}929.
\textbf{11.} Real schedules are made on forecasts; a price-maker's bids also move the prices it
trades. \textbf{12.} Intraday arbitrage, balancing reserves (capacity payments), imbalance
optimisation for a portfolio, and grid services. \textbf{13.} They narrow: more buyers at midday
raise the cheap prices, more sellers in the evening lower the dear ones. \textbf{14.} Each cycle
wears the cells; capping cycles trades revenue today against capacity later. \textbf{15.} More,
shorter price steps to arbitrage, and a better match to the battery's fast response.
\textbf{16.} Both come from the same midday price collapse: what lowers solar's \omterm{def:m3:power-markets-trading:capture}{capture price}
widens the spread batteries trade. \textbf{17.} Other batteries, pumped hydro, flexible industrial
demand, electrolysers and exporters. \textbf{18.} Only while solar keeps outgrowing flexibility;
it erodes as storage is built. \textbf{19.} \emph{Named result:} EUR~74{,}546 per MW in 2025
(EUR~65{,}968 in 2024), one cycle a day, perfect foresight. \textbf{20.} Moving energy in time:
it buys the hours others cannot use and sells the hours they cannot serve.
\end{solution}

\begin{solution}{iq:m3:power-markets-trading:1}
The price at which the grid operator settles a party's deviation from its schedule, derived from
the balancing energy it activated. Being short is cheap when the system is long: the \omterm{def:m3:power-markets-trading:imbalance}{imbalance
price} is then low (even negative), so buying the missing energy from the operator costs little.
\iqlookfor{system direction drives the price; the sign of your error relative to the system's.}
\end{solution}

\begin{solution}{iq:m3:power-markets-trading:2}
Solar produces in the same hours as all other solar, and those hours' prices fall as solar output
rises; the output-weighted price is therefore below the time average. Each new solar park deepens
the midday price fall, so the capture rate keeps falling unless storage or demand absorbs it.
\iqlookfor{self-cannibalisation.}
\end{solution}

\begin{solution}{iq:m3:power-markets-trading:3}
Forecast the hourly output profile (weather, degradation, curtailment) and the hourly price
curve for ten years (fundamental or scenario models); value the output at its \omterm{def:m3:power-markets-trading:capture}{capture price}, not
the baseload; add imbalance and shape costs; discount; compare with the PPA price; stress the
price scenarios and the correlation between output and price.
\iqlookfor{\omterm{def:m3:power-markets-trading:capture}{capture price}, profile and correlation, scenarios.}
\end{solution}

\begin{solution}{iq:m3:power-markets-trading:4}
A financial day-ahead position at a node, settled at the real-time price: it arbitrages the
day-ahead/real-time spread and aligns the two. It can be used to move day-ahead congestion to
benefit a trader's FTRs, losing on the virtual and gaining more on the rights; rules withhold FTR
payments when that pattern appears.
\iqlookfor{convergence role, and the FTR cross-product manipulation.}
\end{solution}

\begin{solution}{iq:m3:power-markets-trading:5}
The error's size and, more, its correlation with the system's imbalance and hence with the
\omterm{def:m3:power-markets-trading:imbalance}{imbalance price}: an error in the system's direction is costly, one against it is rewarded. A
portfolio whose errors correlate with the zone's (all wind forecasts wrong together) pays most.
\iqlookfor{correlation with the system, not only variance.}
\end{solution}

\begin{solution}{iq:m3:power-markets-trading:6}
A stochastic optimisation over price scenarios for the three markets, with the battery's power,
energy, efficiency and cycle constraints, reserve capacity commitments that remove headroom, and
sequential decisions (day-ahead first, intraday and reserves later): a two-stage or rolling
model, solved as a linear programme per scenario set, with degradation cost per cycle.
\iqlookfor{co-optimisation, sequential markets, constraints and uncertainty.}
\end{solution}
